% >>> ARCHIVO GENERADO por scripts/build_figures.py. Se sobrescribe siempre (se carga desde results_macros.tex). <<<
\renewcommand{\figcaption}[1]{\ifcsname figcap@#1\endcsname\csname figcap@#1\endcsname\else\textcolor{red}{\textbf{[??caption \detokenize{#1}]}}\fi}
\expandafter\def\csname figcap@class_distribution\endcsname{Imágenes por clase en el conjunto de entrenamiento.}
\expandafter\def\csname figcap@resolution_sweep\endcsname{Mean recall de validación (continua) y de train (discontinua) según la resolución de entrada, con una capa oculta de 256. Se muestran 10 de 16 runs: todas las resoluciones con lr = 1e-3 y, en las resoluciones con varias lr probadas, la de mejor mean recall de validación.}
\expandafter\def\csname figcap@capacity\endcsname{Accuracy de train (discontinua) y de validación (continua) de las 12 arquitecturas probadas según su número de parámetros, una línea por número de capas ocultas.}
\expandafter\def\csname figcap@lr_sweep\endcsname{Mean recall de validación (continua) y de train (discontinua) según la learning rate, para Adam y SGD con momentum (los 8 runs de la fase 5.1, entrenamientos de hasta 20 épocas).}
\expandafter\def\csname figcap@stopping\endcsname{Punto de parada del modelo final: mean recall de validación por época para cada semilla. La estrella marca la mejor época (la que se conserva) y la cruz, la época en que se detiene EarlyStopping.}
\expandafter\def\csname figcap@progress\endcsname{Mean recall y accuracy de validación tras cada fase (fases 4 y 5: media de 3 semillas).}
\expandafter\def\csname figcap@per_class_recall\endcsname{Recall de validación por clase: baseline frente al modelo final (media de 3 semillas).}
